\section{Introduction}

In February 2018, after decades of globalization and economic integration, the Trump Administration imposed the first of eight major waves of tariffs targeting both China and other trading partners across a variety of goods, initiating the first of two trade wars. These raised the cost of trade for countries exporting to the United States, though numerous studies have found that tariff costs were passed entirely onto US consumers \autocite{fajgelbaumEconomicImpactsUS2022}. Tariffs also both substantially and persistently reduced trade volumes \autocite{amiti_whos_2020}. These findings imply that tariffs are a force of economic disintegration, where prices are higher when we trade, and so we trade less. Yet this trade war arrived in the wake of decades of efforts to deepen global economic ties. This paper asks whether those efforts were in vain. It asks whether the hallmark of integration --- the free trade agreement --- which over the decades has moved beyond a tool to reduce tariffs and increasingly into policy tools to reduce other barriers to trade \autocite{hofmann2017}, helps countries weather the storms trade wars incite. It asks whether deeper trade agreements made countries more resilient to the repeated unilateral tariff shock that was the 2018 trade war. 

The primary consequences of tariffs, the literature finds, were higher prices and less trade. Many studies have found complete pass-through of tariffs to prices paid by importers at the border \autocite{amiti_whos_2020, fajgelbaum2020, cavalloTariffPassThroughBorder2021}, meaning that foreign firms raise prices by exactly the tariff amount, forcing US firms to bear their cost entirely. Interestingly, this is inconsistent with international trade theory, which has long posited that tariffs imposed by a large country would result in incomplete pass-through due to its market power in the global economy \autocite{fajgelbaumEconomicImpactsUS2022}. To retain access to the country's large market, global exporters would reduce their prices, thus absorbing some of the tariff costs. Empirical evidence supported this, consistently finding incomplete pass-through.\footnote{For a more detailed review of tariff pass-through in general and due to the 2018 trade war, see \autocite{fajgelbaumEconomicImpactsUS2022} section 4.1.} Meanwhile, tariffs also reduced trade volumes, defined as the dollar value of imports in a given period \autocite{fajgelbaum2020, amiti2019_tariffs, amiti_whos_2020}. This more closely reflects traditional international trade theory, which predicts tariffs will lower export volumes to the tariffing country. 

Deep trade agreements --- defined as trade agreements that cover more policy areas --- have been shown to lead to higher trade. Deep trade agreements reduce both tariff and non-tariff barriers to trade. The first objective of trade agreements has historically been to reduce tariffs beyond WTO requirements --- tariff reductions are the most common provision in trade agreements, while non-tariff policy provisions vary to a greater degree \autocite{hofmann2017}. Reductions in non-tariff barriers through deep trade agreements have been shown to lead to higher trade creation while causing less trade diversion \autocite{mattooTradeCreationTrade2022}. Greater trade creation, or increased trade between trade agreement members, is a result of discriminatory policies within agreements that affect only those countries who are party to the agreement (e.g. preferential access, mutual recognition, or regulatory alignment). Reduced trade diversion means that deeper agreements reduce the likelihood that agreement members substitute each others' goods for products previously imported from non-members. These are a result of non-discriminatory policies (e.g. competition policy or customs improvements) that apply to all trade relationships. Finally, trade agreements reduce uncertainty \autocite{TPU_literature_review, carballo2022}.\footnote{For further literature on the economic impacts of deep trade agreements, see the literature review by \textcite{fernandes2021}.} These four factors --- reduced tariffs, discriminatory and non-discriminatory policies, and reduced uncertainty --- not only increase trade, but typically reduce costs in the long run. Three of the four factors, excluding only non-discriminatory policies, affect only agreement partners. Under this framework and under conditions of established agreements (where adjustment costs due to factors like regulatory alignment have already been incurred and overcome), I predict that agreement participants have relatively lower trade costs compared to non-participants, creating a wedge between these groups' trade costs. Therefore, in the case of a unilaterally-imposed tariff shock, this wedge helps firms from countries that share bilateral trade agreements with the US to more easily absorb tariff costs, leading to lower pass-through. Meanwhile, I predict that reduced uncertainty and enhanced price competitiveness under these conditions would also lead to a smaller decline in trade volumes.

As such, this paper poses the following question: did the effect of tariffs on countries' trade with the United States differ depending on the depth of countries' trade agreements with the US? In doing so, it asks whether more institutionalized trade relationships are more resilient to trade disruptions as a result of protectionism. It contributes to two main strands of the literature. First, it extends the literature on tariffs' welfare impacts, exploring heterogeneities in tariffs' effects. Understanding how tariffs affect different groups, products, and industries is important to designing effective trade policy. Second, it contributes a new perspective on trade agreements as resilience-enhancing tools, which is all the more important in today's world of geoeconomic instability and unpredictability. As countries like the EU seek to use trade agreements to diversify and institutionalize their trade relationships, understanding trade agreements' diverse impacts can help design agreements more aligned with policy objectives. 

This paper studies whether trade agreements do indeed lead to lower pass-through and lower volume reduction. Extending \textcite{amiti_whos_2020}, I construct an event-study framework with a triple-differences interpretation that identifies the differential effect of trade agreements on prices and volumes.\footnote{For a formalization of the triple-difference estimator, see \textcite{oldenTripleDifferenceEstimator2022}.} I examine the differential effect of agreements on two outcome variables, prices and volumes, examining the first three waves of tariffs. These were imposed through Section 201 and 232 tariffs, applying non-discriminatorily across all countries importing solar panels, washing machines, steel products, and aluminum products. It turns out the US's 20 trade agreements are all primarily deep, meaning there is no sufficient sample of shallow trade agreements to compare their effect to. As such, the paper proceeds with studying the differential effects of trade agreements in general with respect to prices and volumes following exposure to tariffs. 

My results find that countries with bilateral US trade agreements were more resilient to the 2018 tariff shock than countries without. Following tariff implementation, import volumes from agreement countries rose while volumes from non-agreement countries fell, suggesting that institutionalized trade relationships moderated the impact of unilateral tariffs. This aggregate effect is driven by steel products and exhibits a similar but short-lived trend for non-steel products, indicating that the resilience-enhancing role of trade agreements may operate differently depending on product characteristics. The price evidence is consistent with this pattern but cannot be causally identified due to pre-trend violations within product subsamples. These patterns are consistent with the cost-buffer mechanism developed in this paper, in which prior tariff and non-tariff barrier reductions leave agreement-country exporters better positioned to absorb tariff shocks.

The remainder of this paper is organized as follows. Section~\ref{sec:data} describes the data sources and presents descriptive statistics on tariffs, trade agreements, and trade flows. Section~\ref{sec:models} walks through a model to calibrate the restricted sample with the foundational paper, the main triple-difference model, and an examination of heterogeneity between steel and non-steel products due to the high share of the sample constituted by steel products. Section~\ref{sec:results} describes trade agreements' differential price and volume elasticities with respect to tariffs, and Section~\ref{sec:discussion} interprets results, highlighting data issues and assumption violations, along with a discussion of limitations and avenues for future research. Finally, Section~\ref{sec:conclusion} concludes. 

